\documentclass[11pt]{article}

\usepackage[letterpaper,margin=1in]{geometry}
\usepackage{amsmath,amssymb}
\usepackage{booktabs}
\usepackage{longtable}
\usepackage{array}
\usepackage{xcolor}
\usepackage{listings}
\usepackage{url}
\usepackage{hyperref}
\usepackage{titlesec}
\usepackage{parskip}
\usepackage{fancyhdr}
\usepackage{enumitem}

\hypersetup{
  colorlinks=true,
  linkcolor=black,
  urlcolor=blue!50!black,
  citecolor=black,
  pdftitle={Token Importance Unit --- Interface Specification},
  pdfauthor={LonghornSilicon},
  pdfsubject={ISA Specification, Version tiu-isa-0.1}
}

\titleformat{\section}{\normalfont\Large\bfseries}{\thesection.}{0.6em}{}
\titleformat{\subsection}{\normalfont\large\bfseries}{\thesubsection}{0.6em}{}
\titleformat{\subsubsection}{\normalfont\normalsize\bfseries}{\thesubsubsection}{0.6em}{}

\pagestyle{fancy}
\fancyhf{}
\lhead{\small Token Importance Unit ISA}
\rhead{\small \texttt{tiu-isa-0.1}}
\cfoot{\thepage}
\renewcommand{\headrulewidth}{0.4pt}

\definecolor{codebg}{gray}{0.97}
\definecolor{codekey}{rgb}{0.13,0.29,0.53}
\definecolor{codecom}{rgb}{0.25,0.5,0.25}
\lstset{
  basicstyle=\ttfamily\small,
  backgroundcolor=\color{codebg},
  frame=single,
  framerule=0.4pt,
  rulecolor=\color{black!30},
  xleftmargin=0.5em,
  xrightmargin=0.5em,
  aboveskip=0.6em,
  belowskip=0.6em,
  showstringspaces=false,
  breaklines=true,
  keywordstyle=\color{codekey}\bfseries,
  commentstyle=\color{codecom}\itshape,
  numbers=none
}

\setlist{itemsep=1pt,topsep=2pt}

% --- table breathing room ---------------------------------------------------
% Roomier rows, a touch more column padding, and clear vertical separation
% between floats and the surrounding body text (important with parskip, where
% floats otherwise butt against paragraphs).
\renewcommand{\arraystretch}{1.25}
\setlength{\tabcolsep}{7pt}
\setlength{\textfloatsep}{18pt plus 2pt minus 2pt}
\setlength{\intextsep}{16pt plus 3pt minus 2pt}
\setlength{\floatsep}{16pt plus 2pt minus 2pt}

\begin{document}

%==============================================================================
% Title page
%==============================================================================
\thispagestyle{empty}
\begin{center}
\vspace*{1.2in}
{\Huge\bfseries Token Importance Unit}\\[0.4em]
{\LARGE\bfseries Interface Specification}\\[1.2em]

{\large \textsc{LonghornSilicon LLM Inference Accelerator}}\\[0.4em]
{\large Block 3 of 4 --- H2O Heavy-Hitter Eviction Core}\\[2em]

\rule{0.7\textwidth}{0.4pt}\\[1em]

\begin{tabular}{rl}
\textbf{Version}:    & \texttt{tiu-isa-0.1} (first public draft) \\
\textbf{Date}:       & 2026-07-18 \\
\textbf{Status}:     & Pre-tape-out stub for compiler integration \\
\textbf{Author}:     & Chaithu Talasila, The University of Texas at Austin \\
\textbf{Reference}:  & \texttt{LonghornSilicon/token-importance-unit} \\
\end{tabular}\\[1em]
\rule{0.7\textwidth}{0.4pt}

\vfill
\begin{minipage}{0.85\textwidth}
\small
\textbf{Scope.} This document describes the externally-visible interface to the
\texttt{token\_importance\_unit} block as it will appear on the chip, on an FPGA
prototype (Xilinx ZCU102/104), or in the bit-accurate software model. A compiler
backend targeting the TIU writes against this interface; the hardware
implementation must conform to it. This specification is stable for the Token
Importance Unit block only and will be unified with the rest of the
LonghornSilicon ISA when the Attention Compute Unit, KV Cache Engine, and Memory
Hierarchy Controller blocks land.
\end{minipage}

\vspace{0.3in}
\end{center}

\clearpage
\tableofcontents
\clearpage

%==============================================================================
% 1. Block overview
%==============================================================================
\section{Block Overview}
\label{sec:overview}

The Token Importance Unit is the H2O heavy-hitter eviction core for the KV cache.
It maintains \texttt{N\_SLOTS} KV-cache slots, each holding an accumulated
attention-mass score (the heavy-hitter-oracle statistic). It exposes three
operations and one combinational tier read:
\begin{align*}
\texttt{LOAD}(s)      &:\; \texttt{score}[s] \leftarrow 0,\; \texttt{valid}[s] \leftarrow 1 \\
\texttt{ACC}(s, w)    &:\; \texttt{score}[s] \leftarrow \mathrm{sat\_add}(\texttt{score}[s], w) \\
\texttt{EVICT}()      &:\; \rightarrow \arg\min_{k:\,\texttt{valid}[k]} \texttt{score}[k]\;\;\text{(then free it)} \\
\texttt{TIER}(\theta) &:\; \texttt{tier\_keep}[k] = \texttt{valid}[k] \wedge (\texttt{score}[k] \geq \theta)
\end{align*}

\begin{itemize}
\item \textbf{ACC} adds the per-step attention mass a cached token just received,
  one weight per cycle, saturating at $\texttt{SCORE\_MAX} = 2^{\texttt{SCORE\_WIDTH}}-1$.
  Accumulators are \emph{distributed} (one saturating adder per slot, no shared
  broadcast net).
\item \textbf{LOAD} installs a fresh token: zero its score, mark it valid. LOAD
  has priority over ACC when both target the same slot in the same cycle.
\item \textbf{EVICT} runs a \emph{serialized argmin} (one comparator, no wide
  combinational min tree) over the valid slots, returns the minimum-mass slot ---
  the token to drop --- then frees it.
\item \textbf{TIER} is combinational: \texttt{N\_SLOTS} parallel comparators emit
  a per-slot keep/demote flag against the programmable \texttt{TIER\_THRESHOLD}.
\end{itemize}

\medskip
\noindent\textbf{The H2O recent-window and KV-budget bookkeeping is a
control-layer wrapper, not part of this core.} The silicon core is
\emph{accumulate + argmin + tier}; the software layer decides \emph{when} to LOAD
and EVICT (see \path{sw/reference_model/example_compiler_use.py}). This mirrors
how block~1 shipped just the ratio gate and left the tile loop to the compiler.

\medskip
\noindent\textbf{Latency:} \texttt{LOAD} and \texttt{ACC} complete in one cycle
from idle; \texttt{EVICT} takes $\texttt{N\_SLOTS}+1$ cycles (the serialized
argmin scan plus the result). \texttt{TIER} is combinational.\\
\textbf{Throughput:} one \texttt{LOAD}/\texttt{ACC} per cycle; one eviction per
$\texttt{N\_SLOTS}+1$ cycles.\\
\textbf{Footprint:} 95 flip-flops at $\texttt{N\_SLOTS}=8$; $13\,833~\mu$m$^2$
die, $\sim$0.58\,\textmu W at SkyWater Sky130A ($\texttt{N\_SLOTS}=4$ physical
proxy; see \S\ref{sec:syntime} and \path{docs/sky130_signoff.md}).\\
\textbf{Bit-exact reference:} \path{sw/reference_model/tiu_ref.py} (40/40 RTL
replay evictions match).

%==============================================================================
% 2. Memory map
%==============================================================================
\section{Address Space and Memory Map}
\label{sec:memmap}

The block exposes a 256-byte AXI-Lite slave register window. All registers are
32-bit, word-aligned. The base address is set at chip integration time (e.g.,
\texttt{0x4000\_2000} on the ZCU102 PYNQ overlay). $\texttt{SLOT\_WIDTH} =
\lceil\log_2\texttt{N\_SLOTS}\rceil$.

\begin{table}[h]
\centering
\small
\caption{AXI-Lite register window (offsets relative to base address).}
\label{tab:regs}
\begin{tabular}{@{}l l c l p{6.4cm}@{}}
\toprule
Offset & Name & Access & Reset & Purpose \\
\midrule
\texttt{0x00} & \texttt{CTRL}              & RW  & 0x0   & bit[0]: \texttt{soft\_reset} (write-1 to pulse; clears \texttt{valid[]}); bit[1]: \texttt{enable} \\
\texttt{0x04} & \texttt{STATUS}            & R   & 0x1   & bit[0]: \texttt{idle} (= \texttt{!busy}); bit[1]: \texttt{evict\_valid}; bit[2]: \texttt{scan\_in\_progress} (= \texttt{busy}) \\
\texttt{0x08} & \texttt{INFO\_N\_SLOTS}    & R   & (syn) & Synthesis-time \texttt{N\_SLOTS} \\
\texttt{0x0C} & \texttt{INFO\_SCORE\_WIDTH}& R   & (syn) & Bits per accumulated-mass score \\
\texttt{0x10} & \texttt{INFO\_WEIGHT\_WIDTH}& R  & (syn) & Bits per \texttt{ACC} increment \\
\texttt{0x14} & \texttt{INFO\_VERSION}     & R   & 0x0001& bits[15:8] major, bits[7:0] minor \\
\texttt{0x18} & \texttt{TIER\_THRESHOLD}   & RW  & 0x0   & Accumulated-mass keep/demote boundary (\texttt{SCORE\_WIDTH} bits) \\
\texttt{0x1C} & \texttt{OP\_LOAD}          & W   & ---   & Write slot in bits[\texttt{SLOT\_WIDTH}$-$1:0] $\to$ LOAD \\
\texttt{0x20} & \texttt{OP\_ACC}           & W   & ---   & bits[\texttt{SLOT\_WIDTH}$-$1:0]=slot; bits[16$+$]=weight $\to$ ACC \\
\texttt{0x24} & \texttt{OP\_EVICT}         & W   & ---   & Write-1 to bit[0] $\to$ trigger a serialized-argmin scan \\
\texttt{0x28} & \texttt{EVICT\_RESULT}     & R   & 0x0   & bit[31]: \texttt{valid}; bits[\texttt{SLOT\_WIDTH}$-$1:0]: victim slot \\
\texttt{0x2C} & \texttt{TIER\_KEEP}        & R   & 0x0   & bit[$k$] = \texttt{tier\_keep}[$k$] (1 = keep/CQ-8, 0 = demote/CQ-4) \\
\texttt{0x30} & \texttt{SLOT\_VALID}       & R   & 0x0   & bit[$k$] = \texttt{valid}[$k$] (diagnostic occupancy) \\
\texttt{0x34} & \texttt{SCORE\_SEL}        & RW  & 0x0   & Slot index selecting \texttt{SCORE\_READ} (diagnostic) \\
\texttt{0x38} & \texttt{SCORE\_READ}       & R   & ---   & Accumulated mass of the \texttt{SCORE\_SEL} slot (diagnostic) \\
\bottomrule
\end{tabular}
\end{table}

\noindent\textbf{Conventions.} \texttt{RW} = read/write; \texttt{R} = read-only;
\texttt{W} = write-only trigger; (syn) = value fixed at synthesis time. Writes to
read-only registers are silently dropped; reading reserved offsets returns zero.
\texttt{OP\_*} writes are ignored while \texttt{STATUS.scan\_in\_progress} is
asserted; a compiler polls \texttt{STATUS.idle} before the next op, or uses the
streaming interface (\S\ref{sec:stream}), which handles backpressure in hardware.

%==============================================================================
% 3. Streaming interface
%==============================================================================
\section{Streaming Op Interface}
\label{sec:stream}

For high throughput the op stream is also exposed as an AXI-Stream channel; the
AXI-Lite window is then used only for \texttt{INFO\_*}, \texttt{TIER\_THRESHOLD},
and diagnostics. This is the recommended path --- one op per cycle with hardware
backpressure, matching the RTL's native op ports.

\subsection{\texttt{s\_axis\_ops} --- Operation Input (Slave)}

\begin{table}[h]
\centering
\small
\begin{tabular}{@{}l c c p{7.2cm}@{}}
\toprule
Signal & Width & Direction & Purpose \\
\midrule
\texttt{tdata}  & $2{+}\texttt{SLOT\_WIDTH}{+}\texttt{WEIGHT\_WIDTH}$ & in & \{\texttt{op}[1:0], \texttt{slot}, \texttt{weight}\} \\
\texttt{tvalid} & 1 & in  & Handshake: op is valid \\
\texttt{tready} & 1 & out & Handshake: block is ready to accept \\
\bottomrule
\end{tabular}
\end{table}

\noindent\texttt{op} encoding: $0=\texttt{ACC}$, $1=\texttt{LOAD}$,
$2=\texttt{EVICT}$ (matching the golden trace and
\path{analysis/gen_tiu_testvectors.py}). For \texttt{LOAD}, \texttt{weight} is
ignored; for \texttt{EVICT}, \texttt{slot} and \texttt{weight} are ignored.
\texttt{tready} deasserts for \texttt{N\_SLOTS} cycles after an \texttt{EVICT}
while the argmin scan runs.

\subsection{\texttt{m\_axis\_evict} --- Eviction Output (Master)}

\begin{table}[h]
\centering
\small
\begin{tabular}{@{}l c c p{7.2cm}@{}}
\toprule
Signal & Width & Direction & Purpose \\
\midrule
\texttt{tdata}  & \texttt{SLOT\_WIDTH} & out & The evicted (minimum-mass) victim slot \\
\texttt{tvalid} & 1 & out & Handshake: victim is valid (one beat per \texttt{EVICT}) \\
\texttt{tready} & 1 & in  & Handshake: downstream (KVCE) is ready \\
\bottomrule
\end{tabular}
\end{table}

\noindent One beat per \texttt{EVICT}, emitted $\texttt{N\_SLOTS}+1$ cycles after
the \texttt{EVICT} is accepted. \texttt{tier\_keep} is not streamed --- it is a
continuously-valid combinational output the KV cache engine samples when it
(re)compresses a token's value (\S\ref{sec:tier}).

%==============================================================================
% 4. Logical operations
%==============================================================================
\section{Logical Operations}
\label{sec:ops}

\begin{table}[h]
\centering
\small
\begin{tabular}{@{}l l l p{5.6cm}@{}}
\toprule
Operation           & Inputs        & Outputs      & Description \\
\midrule
\texttt{TIU\_QUERY}     & ---           & INFO struct  & Read all \texttt{INFO\_*} registers in one transaction \\
\texttt{TIU\_RESET}     & ---           & ---          & Soft reset: clear \texttt{valid[]} (empty the cache) \\
\texttt{TIU\_ENABLE}    & ---           & ---          & Set bit[1] of \texttt{CTRL} \\
\texttt{TIU\_LOAD}      & slot          & ---          & Install a fresh token (score $:=0$, valid $:=1$) \\
\texttt{TIU\_ACC}       & slot, weight  & ---          & Add \texttt{weight} of mass to \texttt{slot} (saturating) \\
\texttt{TIU\_EVICT}     & ---           & victim slot  & Run argmin; return + free the minimum-mass valid slot \\
\texttt{TIU\_SET\_TIER} & threshold     & ---          & Program \texttt{TIER\_THRESHOLD} \\
\texttt{TIU\_READ\_TIER}& ---           & keep bitmap  & Read \texttt{TIER\_KEEP} (per-slot keep/demote) \\
\bottomrule
\end{tabular}
\end{table}

\noindent The block has no other externally-observable side effects: no
clock-gating register, no power state machine, and no DFT scan-chain control at
this stub level (those land on the chip-level ISA).

\subsection{Worked Lowering Example}
\label{sec:lowering}

How a representative compiler IR operation lowers to a sequence of TIU ops. The
example uses a generic \texttt{kv\_cache\_append} op any silicon-agnostic
compiler is likely to have; the H2O wrapper is compiler-side.

\begin{lstlisting}[basicstyle=\ttfamily\small]
info = TIU_QUERY()                 // -> n_slots=8, score_width=8, weight_width=8
TIU_RESET(); TIU_ENABLE()
C = min(round(0.25 * ctx_len), info.n_slots)   // gold 25% KV budget

for t in tokens:
    free = first_invalid_slot()
    if free is None:               // cache full -> evict the argmin victim
        victim = TIU_EVICT()       // N_SLOTS+1 cycles, returns slot
        kvce_drop(slot=victim)     // KVCE frees this token's K and V
        free = victim
    TIU_LOAD(free)                 // install token t
    for (slot, w) in attn_mass_over_cache(t):
        TIU_ACC(slot, w)           // 1 cycle each; accumulate received mass
\end{lstlisting}

\noindent\texttt{attn\_mass\_over\_cache(t)} is evaluated \emph{after} the
\texttt{TIU\_LOAD}, keyed by each slot's \emph{current} occupant (token $t$'s
self-attention included). A wrapper that instead precomputes weights against the
pre-step cache must issue those \texttt{TIU\_ACC}s \emph{before} the
\texttt{TIU\_EVICT} --- otherwise the evicted token's mass is credited to the
fresh token reusing its slot (cf.\ \texttt{H2OScheduler.step} in the reference
model).

\textbf{Value-tier pass} (when the KVCE (re)compresses a surviving token's value):
\begin{lstlisting}[basicstyle=\ttfamily\small]
TIU_SET_TIER(threshold)            // e.g. a percentile of accumulated mass
keep = TIU_READ_TIER()             // bit[s] = 1 -> value @ CQ-8, else CQ-4
for s in occupied_slots:
    kvce_store_value(slot=s, bits = 8 if keep[s] else 4)
\end{lstlisting}

\noindent The same lowering works whether the compiler's target is the Python
reference model (phase~0), the FPGA prototype (phase~1), or the silicon chip
(phase~3). Only the runtime implementation of the \texttt{TIU\_*} operations
changes underneath --- Python call, AXI register write, or PCIe MMIO.

%==============================================================================
% 5. Tier handshake
%==============================================================================
\section{Tier Handshake to the KV Cache Engine}
\label{sec:tier}

The TIU tells the KV Cache Engine (block~2) how to treat each cached token ---
\textbf{keep}, \textbf{demote}, or \textbf{evict} --- via two channels:

\begin{enumerate}
\item \textbf{Evict} (\texttt{m\_axis\_evict} / \texttt{EVICT\_RESULT}): the
  minimum-mass victim's K and V are \emph{dropped} and the slot freed. This is
  the per-token lever that applies to \emph{both} K and V.
\item \textbf{Keep / demote} (\texttt{TIER\_KEEP} against
  \texttt{TIER\_THRESHOLD}): for a surviving token in slot $s$, the KVCE reads
  \texttt{tier\_keep}[$s$] when it (re)compresses that token's \emph{value}:
  \begin{itemize}
  \item \texttt{tier\_keep}[$s$]$=1$ (heavy hitter) $\to$ store the \textbf{value
    at CQ-8} (per-token INT8).
  \item \texttt{tier\_keep}[$s$]$=0$ (demote) $\to$ store the \textbf{value at
    CQ-4} (per-token INT4).
  \end{itemize}
\end{enumerate}

\noindent The tier is a per-token \textbf{value}-precision lever only. Keys stay
uniform per-channel (per-token key demotion collapses GQA; measured $-0.17$ in
the full stack), so the whole key cache shares one ChannelQuant key tier, set
globally. \texttt{tier\_keep} is emitted as \texttt{N\_SLOTS} parallel
comparators (no read-mux), so it adds no fanout to the argmin datapath --- this is
what kept the Sky130 sign-off at 0 violations after the port was added. The full
protocol, verified end-to-end with the ACU precision controller in the loop, is in
\path{docs/tier_handshake.md}.

%==============================================================================
% 6. Synthesis-time configuration
%==============================================================================
\section{Synthesis-Time Configuration}
\label{sec:syntime}

The following parameters are baked in at synthesis and cannot be changed at
runtime. A compiler reads them via the \texttt{INFO\_*} registers and tailors its
code generation accordingly.

\begin{table}[h]
\centering
\small
\begin{tabular}{@{}l c l p{5.5cm}@{}}
\toprule
Parameter          & Default & Range (validated)      & Notes \\
\midrule
\texttt{N\_SLOTS}     & 8 & 2, 4, 8, 16, 32   & KV-cache slots (per head, or per shared group) \\
\texttt{SCORE\_WIDTH} & 8 & 6, 8, 10, 12, 16  & Accumulated-mass width; \textbf{8 is loss-free} \\
\texttt{WEIGHT\_WIDTH}& 8 & 4, 6, 8           & Per-step attention-mass increment width \\
\bottomrule
\end{tabular}
\end{table}

\noindent \texttt{SCORE\_WIDTH} $=8$ is set from the accumulator-bit-width study
(\path{docs/findings/h2o-deep-analysis.md}): 8 bits is loss-free for the eviction
ranking on Qwen2 ($\Delta$ $-0.002$ acc\_norm); 6 bits breaks ($-0.022$). Fewer
score bits also shrink the argmin read-mux, which drove the Sky130 max-transition
closure.

\medskip
\noindent\textbf{Physical proxy: \texttt{N\_SLOTS} $=4$.} The shipped RTL default
is \texttt{N\_SLOTS} $=8$ (what the testbench verifies bit-exact). The Sky130
physical run synthesizes an \texttt{N\_SLOTS} $=4$ proxy via
\texttt{SYNTH\_PARAMETERS} --- exactly as the KV Cache Engine ships a
\texttt{VECTOR\_DIM} $=8$ proxy --- to shrink the argmin-mux fanout and clock
tree so the clock-root fanout clears the flow's limit; the datapath is
parameter-identical. FF count: \textbf{95 @ N\_SLOTS=8} (real),
\textbf{53 @ N\_SLOTS=4} (proxy). Real \texttt{N\_SLOTS} is set
per-instantiation. The count tracks the closed form
\[
  \mathrm{FFs} = \texttt{N\_SLOTS}\cdot(\texttt{SCORE\_WIDTH}{+}1) + \texttt{SCORE\_WIDTH} + 3\cdot\texttt{SLOT\_WIDTH} + 3
\]
(92 @ N\_SLOTS=8; yosys keeps $\sim{+}3$ slot-index FFs un-merged $\to$ 95
synthesized, the value the CI gate pins).

%==============================================================================
% 7. Bit-accurate reference
%==============================================================================
\section{Bit-Accurate Reference}
\label{sec:ref}

The Python model at \path{sw/reference_model/tiu_ref.py} is the canonical
bit-accurate reference for this ISA. A compiler can verify its codegen against
the model directly:

\begin{lstlisting}[language=Python]
from tiu_ref import TokenImportanceUnit

tiu = TokenImportanceUnit()
tiu.load(0); tiu.acc(0, 200)
victim = tiu.evict()             # exactly what the chip will return
keep   = tiu.tier(threshold=40)  # per-slot keep/demote, exactly as tier_keep
\end{lstlisting}

\noindent The model is verified bit-exact against the RTL by replaying the
committed golden trace \path{rtl/tb/testvectors/tiu_trace.hex} (48 real Qwen2
tokens, 40 evictions) through both the model and the SystemVerilog replay
testbench: all 40 eviction victims agree
(\path{sw/reference_model/test_tiu_ref.py}). Any divergence between the model and
the chip is a bug in this specification, not in the chip.

%==============================================================================
% 8. Integration phases
%==============================================================================
\section{Integration Phases}
\label{sec:phases}

\begin{table}[ht]
\centering
\small
\renewcommand{\arraystretch}{1.25}
\begin{tabular}{@{}c p{3.0cm} p{4.0cm} p{5.0cm}@{}}
\toprule
Phase & Timeline & Compiler targets & We provide \\
\midrule
0 & now & Bit-accurate Python reference model & Model + ISA + parity test \\
1 & when ZCU102/104 arrives & AXI-Lite + AXI-Stream on Zynq UltraScale+ & Vivado bitstream + PYNQ driver \\
2 & after all four blocks land & Same AXI interface; tier handshake to KVCE on hardware & Multi-block FPGA project \\
3 & post-tape-out (2027+) & PCIe-attached custom chip & TSMC 16FFC silicon + Linux driver \\
\bottomrule
\end{tabular}
\end{table}

\noindent Throughout all phases, the interface described in this document is the
stable contract.

%==============================================================================
% 9. Change log
%==============================================================================
\section{Change Log}
\label{sec:changelog}

\begin{itemize}
\item \texttt{tiu-isa-0.1} (2026-07-18): First public draft. Stable for the Token
      Importance Unit block only. Documents LOAD/ACC/EVICT plus the tier read,
      the AXI-Lite register and streaming map, the \texttt{N\_SLOTS}$=4$ physical
      proxy, and the KVCE tier handshake.
\end{itemize}

\vspace{0.5in}
\noindent\rule{\textwidth}{0.4pt}\\
\noindent\small\textit{LonghornSilicon Token Importance Unit --- Interface
Specification. This document is open and may be redistributed. The TSMC 16FFC PDK
referenced in the integration phases is NDA-protected and is not part of this
document or the public repository.}

\end{document}
